\Lecturer{Hui Jin}
\Contributor{}
\LectureNumber{08}
\LectureDate{DATE}
\LectureTitle{Turbo Codes Decoding and Design}

\lstset{style=mystyle}

	\MakeScribeTop

%#############################################################
%#############################################################
%#############################################################
%#############################################################

\section{Introduction}
The key innovation of turbo codes is how they use the likelihood data to reconcile differences between the two decoders. Each of the two convolutional decoders generates a hypothesis (with derived likelihoods) for the pattern of m bits in the payload sub-block. The hypothesis bit-patterns are compared, and if they differ, the decoders exchange the derived likelihoods they have for each bit in the hypotheses. Each decoder incorporates the derived likelihood estimates from the other decoder to generate a new hypothesis for the bits in the payload. Then they compare these new hypotheses. This iterative process continues until the two decoders come up with the same hypothesis for the m-bit pattern of the payload, typically in a couple dozen cycles.

\section{Maximum A Posteriori Decoder of Convolutional Codes}
In the lecture of Convolutional Code decoding, we described Maximum Likelihood Decoder, where the decoding decision is to choose the code \(c\) that maximizes the probability \(p(y|c)\). It corresponds to  
\begin{align}
    c_{ML} &= \argmax_{c} \log p(y|c) \nonumber \\ 
           &= \argmax_{c=c_1, c_2, \cdots, c_n} \log \prod_{i=1}^{n}p(y_i|c_i) \nonumber \\
           &= \argmax_{c=c_1, c_2, \cdots, c_n} \sum_{i=1}^{n} \log p(y_i|c_i) \label{eq:ML_Viterbi}
\end{align}
The maximum-likelihood decoder finds the whole sequence of information bits, thus the codeword, that maximizes the received sequence. The decoder admits a linear complexity Viterbi decoding algorithm because of the Markov tree property of the trellis structure. 

For the turbo codes with a concatenated convolution code, the maximum-likelihood decoder cannot be achieved with linear decoding complexity. Rather, the genius of the turbo code inventors lies at the iterative decoding method. The decoding algorithm treats the two convolution codes separately and decodes each one with the help of the other one. They showed with just a few iterative steps, the decoding performance is spectacular. 


 The iterative decoding algorithm works in successive fashion. We decode one convolution code first, obtaining the Maximum A Posteriori (MAP) soft decision of the information bits; then we use such MAP decision as prior, decoding the other convolution code, obtaining the MAP decision from that code. This new MAP decision is subsequently used to decode the first convolution code, starting the second decoding cycle. 

The fundamental decoding block is to obtain the MAP soft decision of the convolution code, given prior information. The algorithm is known as the BCJR algorithm, derived by Bahl, Cocke, Jelinek and Raviv in 1974. The BCJR algorithm is also known as the Forward-Backward algorithm. 

The BCJR algorithm gives the MAP decision on the individual information bits

\section{Turbo Decoding - BCJR Algorithm}

\subsection{A Priori and A Posteriori}
Let us consider a block or convolutional encoder described by a trellis, and a sequence
\(x = x_1x_2 ...x_N\) of \(N\) \(n\)-bit codewords or symbols at its output, where \(x_k\) is the symbol generated by the
encoder at time \(k\). The corresponding information or message input bit, \(uk\), can take on the values \(0\) or \(1\) with an a priori probability \(P(u_k)\), from which we can define the so-called log-likelihood ratio (LLR):

\begin{equation}
    L(u_k) = \ln \frac{P(u_k=1)}{P(u_k=0)}
\end{equation}

Suppose the coded sequence \(x\) is transmitted over a memoryless additive white gaussian noise
(AWGN)channel and  received as the sequence of \(nN\) real numbers \(y=y_1y_2...y_N\).
This sequence is delivered to the decoder and used by the decoder algorithm in order to estimate the original bit sequence \(u_k\). For which the algorithm computes the a posteriori log-likelihood ratio
\(L(u_k| y)\) , a real number defined by the ratio:
\begin{equation}
    L(u_k|y) = \ln \frac{P(u_k=1|y)}{P(u_k=0|y)}
\end{equation}

\subsection{MAP derivation}
We want to calculate \(L(u_k|y) = \ln \frac{P(u_k = 1|y)}{P(u_k = 0|y)}\). If \(L(u_k |y) > 0\), the decoder can make a hard decision that the bit
\(u_k = 1\) was sent, otherwise it will estimate \(u_k = 0\) was sent. Using our simple exemplary convolution code, the trellis has four states; from times k-1 and k there are \(8\) transitions between states \(S_{k-1} = s’\) and next states \(S_k = s\): four result from a \(0\) input bit (those drawn with a thick line) and the other four result from a \(1\) input bit. Each one of these transitions is due unequivocally to a known bit (just notice in the trellis if the branch line is solid or dashed). Therefore, the probabilities that \(u_k = 0\) or \(u_k = 1\), given our knowledge of the sequence \(y\), are equal to the probabilities that the transitions (s’,s) between states correspond to a solid or a dashed line, respectively. Since transitions are mutually exlusive – just one can occur at a given time – the probability that any one of them occurs is equal to the sum of the individual probabilities, that is,
\begin{align*}
    P(u_k=1|y) &= \sum_{s', s, u_k=1}P(s', s|y),\\
    P(u_k=0|y) &= \sum_{s', s, u_k=0}P(s', s|y).  
\end{align*}


\section{Turbo Decoding as an instance of Belief Propagation}


